% The user/kernel protection boundary and the ways to cross it.
% Usage: % The user/kernel protection boundary and the ways to cross it.
% Usage: % The user/kernel protection boundary and the ways to cross it.
% Usage: % The user/kernel protection boundary and the ways to cross it.
% Usage: \input{figures/l01-user-kernel}   (width ~112 mm)
\begin{tikzpicture}[x=1mm, y=1mm, font=\footnotesize\sffamily,
  app/.style={draw, rounded corners=2pt, fill=white, minimum height=8mm,
              minimum width=30mm, inner sep=1pt},
  big/.style={draw, rounded corners=2pt, minimum height=8mm,
              minimum width=112mm, inner sep=1pt},
  lab/.style={font=\scriptsize\sffamily, align=center},
  lvl/.style={font=\scriptsize\sffamily, text=ln-muted, align=right},
  >={Stealth[length=1.8mm]}]
  \node[app] (a) at (28,42) {\iflangja{アプリケーション}{application}};
  \node[app] (b) at (78,42) {\iflangja{アプリケーション}{application}};
  \draw[dashed, ln-rule, thick] (0,26) -- (112,26);
  \node[lvl, anchor=south east] at (112,26.8) {\iflangja{ユーザレベル（非特権）}{user level (unprivileged)}};
  \node[lvl, anchor=north east] at (112,25.2) {\iflangja{カーネルレベル（特権）}{kernel level (privileged)}};
  \node[big, fill=ln-box, thick, draw=ln-accent] (k) at (56,15)
    {\iflangja{OS カーネル}{OS kernel}};
  \node[big, fill=white] (h) at (56,0) {\iflangja{ハードウェア}{hardware}};
  % system call and signal
  \draw[->, thick, ln-accent] (22,38) -- (22,19);
  \node[lab, anchor=east] at (21,32) {\iflangja{システムコール}{system call}};
  \draw[->, thick, ln-accent] (34,19) -- (34,38);
  \node[lab, anchor=west] at (35,32) {\iflangja{シグナル}{signal}};
  % trap
  \draw[->, thick, ln-caution] (78,38) -- (78,19);
  \node[lab, anchor=east] at (77,32)
    {\iflangja{トラップ\\（特権操作）}{trap (privileged\\operation)}};
  % privileged access
  \draw[<->, thick] (56,11) -- (56,4);
  \node[lab, anchor=west] at (57,7.5)
    {\iflangja{ハードウェアへの直接アクセス}{direct hardware access}};
\end{tikzpicture}
   (width ~112 mm)
\begin{tikzpicture}[x=1mm, y=1mm, font=\footnotesize\sffamily,
  app/.style={draw, rounded corners=2pt, fill=white, minimum height=8mm,
              minimum width=30mm, inner sep=1pt},
  big/.style={draw, rounded corners=2pt, minimum height=8mm,
              minimum width=112mm, inner sep=1pt},
  lab/.style={font=\scriptsize\sffamily, align=center},
  lvl/.style={font=\scriptsize\sffamily, text=ln-muted, align=right},
  >={Stealth[length=1.8mm]}]
  \node[app] (a) at (28,42) {\iflangja{アプリケーション}{application}};
  \node[app] (b) at (78,42) {\iflangja{アプリケーション}{application}};
  \draw[dashed, ln-rule, thick] (0,26) -- (112,26);
  \node[lvl, anchor=south east] at (112,26.8) {\iflangja{ユーザレベル（非特権）}{user level (unprivileged)}};
  \node[lvl, anchor=north east] at (112,25.2) {\iflangja{カーネルレベル（特権）}{kernel level (privileged)}};
  \node[big, fill=ln-box, thick, draw=ln-accent] (k) at (56,15)
    {\iflangja{OS カーネル}{OS kernel}};
  \node[big, fill=white] (h) at (56,0) {\iflangja{ハードウェア}{hardware}};
  % system call and signal
  \draw[->, thick, ln-accent] (22,38) -- (22,19);
  \node[lab, anchor=east] at (21,32) {\iflangja{システムコール}{system call}};
  \draw[->, thick, ln-accent] (34,19) -- (34,38);
  \node[lab, anchor=west] at (35,32) {\iflangja{シグナル}{signal}};
  % trap
  \draw[->, thick, ln-caution] (78,38) -- (78,19);
  \node[lab, anchor=east] at (77,32)
    {\iflangja{トラップ\\（特権操作）}{trap (privileged\\operation)}};
  % privileged access
  \draw[<->, thick] (56,11) -- (56,4);
  \node[lab, anchor=west] at (57,7.5)
    {\iflangja{ハードウェアへの直接アクセス}{direct hardware access}};
\end{tikzpicture}
   (width ~112 mm)
\begin{tikzpicture}[x=1mm, y=1mm, font=\footnotesize\sffamily,
  app/.style={draw, rounded corners=2pt, fill=white, minimum height=8mm,
              minimum width=30mm, inner sep=1pt},
  big/.style={draw, rounded corners=2pt, minimum height=8mm,
              minimum width=112mm, inner sep=1pt},
  lab/.style={font=\scriptsize\sffamily, align=center},
  lvl/.style={font=\scriptsize\sffamily, text=ln-muted, align=right},
  >={Stealth[length=1.8mm]}]
  \node[app] (a) at (28,42) {\iflangja{アプリケーション}{application}};
  \node[app] (b) at (78,42) {\iflangja{アプリケーション}{application}};
  \draw[dashed, ln-rule, thick] (0,26) -- (112,26);
  \node[lvl, anchor=south east] at (112,26.8) {\iflangja{ユーザレベル（非特権）}{user level (unprivileged)}};
  \node[lvl, anchor=north east] at (112,25.2) {\iflangja{カーネルレベル（特権）}{kernel level (privileged)}};
  \node[big, fill=ln-box, thick, draw=ln-accent] (k) at (56,15)
    {\iflangja{OS カーネル}{OS kernel}};
  \node[big, fill=white] (h) at (56,0) {\iflangja{ハードウェア}{hardware}};
  % system call and signal
  \draw[->, thick, ln-accent] (22,38) -- (22,19);
  \node[lab, anchor=east] at (21,32) {\iflangja{システムコール}{system call}};
  \draw[->, thick, ln-accent] (34,19) -- (34,38);
  \node[lab, anchor=west] at (35,32) {\iflangja{シグナル}{signal}};
  % trap
  \draw[->, thick, ln-caution] (78,38) -- (78,19);
  \node[lab, anchor=east] at (77,32)
    {\iflangja{トラップ\\（特権操作）}{trap (privileged\\operation)}};
  % privileged access
  \draw[<->, thick] (56,11) -- (56,4);
  \node[lab, anchor=west] at (57,7.5)
    {\iflangja{ハードウェアへの直接アクセス}{direct hardware access}};
\end{tikzpicture}
   (width ~112 mm)
\begin{tikzpicture}[x=1mm, y=1mm, font=\footnotesize\sffamily,
  app/.style={draw, rounded corners=2pt, fill=white, minimum height=8mm,
              minimum width=30mm, inner sep=1pt},
  big/.style={draw, rounded corners=2pt, minimum height=8mm,
              minimum width=112mm, inner sep=1pt},
  lab/.style={font=\scriptsize\sffamily, align=center},
  lvl/.style={font=\scriptsize\sffamily, text=ln-muted, align=right},
  >={Stealth[length=1.8mm]}]
  \node[app] (a) at (28,42) {\iflangja{アプリケーション}{application}};
  \node[app] (b) at (78,42) {\iflangja{アプリケーション}{application}};
  \draw[dashed, ln-rule, thick] (0,26) -- (112,26);
  \node[lvl, anchor=south east] at (112,26.8) {\iflangja{ユーザレベル（非特権）}{user level (unprivileged)}};
  \node[lvl, anchor=north east] at (112,25.2) {\iflangja{カーネルレベル（特権）}{kernel level (privileged)}};
  \node[big, fill=ln-box, thick, draw=ln-accent] (k) at (56,15)
    {\iflangja{OS カーネル}{OS kernel}};
  \node[big, fill=white] (h) at (56,0) {\iflangja{ハードウェア}{hardware}};
  % system call and signal
  \draw[->, thick, ln-accent] (22,38) -- (22,19);
  \node[lab, anchor=east] at (21,32) {\iflangja{システムコール}{system call}};
  \draw[->, thick, ln-accent] (34,19) -- (34,38);
  \node[lab, anchor=west] at (35,32) {\iflangja{シグナル}{signal}};
  % trap
  \draw[->, thick, ln-caution] (78,38) -- (78,19);
  \node[lab, anchor=east] at (77,32)
    {\iflangja{トラップ\\（特権操作）}{trap (privileged\\operation)}};
  % privileged access
  \draw[<->, thick] (56,11) -- (56,4);
  \node[lab, anchor=west] at (57,7.5)
    {\iflangja{ハードウェアへの直接アクセス}{direct hardware access}};
\end{tikzpicture}
